% SOURCE: docs/PAPER_DRAFT_INTRO_DIAGNOSIS.md §1 + integration repairs
% (guarantee-iff phrasing, degree-0 homogeneity, Buffalo "recent",
% [[INSURANCE-PAYOUT-RANGE]] filled, harvest contribution updated). Single-touch.
\section{Introduction}
\label{sec:intro}

The same-seed paired comparison is widely used to evaluate deep-learning
methods: to attribute a small performance difference to a method rather
than run-to-run noise, one launches method and baseline from an identical
seed, holds everything else fixed, and reads the gap as the method's
effect. The instrument looks self-verifying: same seed, same
initialization, same data order, so whatever diverges must be the method.
This paper studies when that intended experimental control fails. A
mid-training probe that advances a stream later consumed by training can
change subsequent stochastic inputs and updates. The fixed-offset account
holds on the aligned prefix; later reconvergence is possible.
State-pure RNG isolation and deterministic execution restore the sufficient
pairing condition, with per-path repeat checks testing its implementation.
The first failure is not automatically a marginal mean bias: it is loss of
the common-random-number coupling and its covariance-based variance reduction.
If the shifted stream preserves the arm's marginal law, more seeds still
estimate the marginal method effect; whether paired variance increases depends
on the change in covariance. Directly selecting the largest reported gap has
nonnegative expected reporting optimism relative to the selected placement's
own marginal mean (Proposition~\ref{prop:harvest}), without iid or Gaussian
assumptions. The observed uplift above the placement average also includes
systematic placement effects, which we separate analytically. Whether this
translates into a positive gain on an untouched test set after validation-based
selection remains unresolved: both primary E1 intervals include zero
(Section~\ref{sec:harvest}).
We read twenty-eight official and field-standard codebases spanning
2018--2026, from noisy-label vision methods to LLM pretraining, and found the
exposure in twenty-two, among them all eleven in the neighborhood the diagnosis
came from, four of which we confirmed broken by re-execution; the selection
step is the machinery those
same pipelines and reporting practices already run, from best-of-$m$ placement
and configuration choices to validation-guarded champion selection.

Our main finding is the \emph{pairing-divergence artifact}. A mid-training
probe that consumes a stream later used by training can change batch orders
or dropout masks and break exact coupling (Proposition~\ref{prop:pairing}). An in-house paired
comparison reporting $+2.97$ fell to $+0.16$ once the probe's RNG was isolated,
and the guarded endpoint motivates the conditional no-op result
(Section~\ref{sec:receipt}); its historical metric zeros are not universal
bitwise certificates.

\paragraph{Contributions.}
(1) We diagnose a failure of same-seed evaluation, formalize a sufficient
pairing condition and the onset of divergence,
verify the tested boundary inside the
HuggingFace Trainer (Figure~\ref{fig:hf-boundary}), audit exposure across
twenty-eight codebases (Section~\ref{sec:prevalence}: eleven core repositories
double-read, twelve field-standard training loops, five flagship
stacks). Separately, selection on the reported quantity yields an observed
$+4.01$ uplift over the placement average in released code. We use
$c_6\bar\sigma$ as a conditional iid-Gaussian benchmark, with same-sample
calibration on one codebase, one dataset and a single rig
(Proposition~\ref{prop:harvest}; Section~\ref{sec:harvest}). A separate nanoGPT
experiment freezes the prediction before evaluating held-out seeds (16 seeds
per split), assessing mean calibration within one additional setting rather
than cross-setting transferability. Its discrepancy interval includes zero.
This held-out-seed check is
distinct from the validation-to-test bridge above. Monitoring and full
$120$-epoch witnesses extend the pairing diagnosis.
(2) We give a practical repair: isolate every probe, pin kernels
(nondeterministic ones reproduce the failure with no RNG sharing at all), and
require an exact repeat floor. On top of it we prove a bitwise no-op
certificate (Corollary~\ref{cor:noop}) by which a repaired selection endpoint
can prove that a comparison it declined to change is unpolluted
(Section~\ref{sec:receipt}). We claim no unconditional accuracy gain for that
endpoint anywhere in this paper; its payout behavior is reported in the appendix.
\contribdiagtwo{}

Common-random-number alignment and selection optimism have established
theoretical foundations; recent work treats event alignment in stochastic
simulation \citep{buffalo2026realizing}. Our contribution is their concrete
connection in deep-learning evaluation: source-level exposure audits,
controlled first-consumer traces, repaired comparisons against probe-free
baselines, and a quantitative assessment of placement selection. The general
selection inequalities clarify the estimand; the released-code diagnosis and
its verification supply the empirical contribution.

The diagnosis is constructive. The repair has three requirements with
different costs: isolation is often a small code change, whereas deterministic
kernels and per-path repeat/state checks add runtime, storage, and sometimes
platform constraints. Under verified premises, a selection endpoint that
serves its raw candidate yields the conditional no-op certificate in
Corollary~\ref{cor:noop}. Historical tuned-regime zero metrics are reported
separately from this certificate: cells without an independent floor are
excluded from certificate counts, and metric equality is not upgraded to
weight or full-state identity. Section~\ref{sec:receipt} and
Figure~\ref{fig:ledger} distinguish the available evidence.
Code, data and analysis are released at
\url{https://github.com/dumingyu2004/silent-unpairing-tmlr2026};
Appendix~\ref{sec:availability} describes the release's contents
and the remaining limits of historical reproduction.
